\documentclass[11pt]{article}
\usepackage[T1]{fontenc}
\usepackage{lmodern,amsmath,amssymb,amsthm}
\usepackage[a4paper,margin=27mm]{geometry}
\usepackage[hidelinks]{hyperref}
\newtheorem{theorem}{Theorem}
\newtheorem{lemma}[theorem]{Lemma}
\newtheorem{corollary}[theorem]{Corollary}
\newcommand{\interior}{\operatorname{int}}
\newcommand{\conv}{\operatorname{conv}}
\newcommand{\area}{\operatorname{area}}
\newcommand{\Rot}{\operatorname{Rot}}
\title{A sector obstruction to two-orbit rotational Stein tilings}
\author{MATH-003 research draft --- LLM-generated}
\date{21 October 2026}
\begin{document}
\maketitle

\noindent\textbf{Status.} This is a proof draft, not an independently reviewed
theorem or an enabled software filter. The local deductions have been checked
within this research iteration. Published premises are identified below;
neither their full proofs nor literature novelty are certified here. The
earlier source is \texttt{rotational\_stein\_structural\_note.tex}.

\section{Question and hypotheses}

Let $D=\{x\in\mathbb R^2:|x|\leq1\}$ and $O=0$. A tiling is a finite
family of congruent compact Jordan regions with pairwise disjoint interiors
whose union is $D$. Reflected congruences are permitted. Assume $N>1$, one tile
$C$ satisfies $O\in\interior C$, and rotation $R=\Rot_O(2\pi/m)$, $m\geq2$,
permutes the geometric tiles. No smoothness, rectifiable tile boundaries,
connected internal contacts, prescribed graph, or maximality of $m$ is assumed.
For a different generator of this same cyclic group, use a suitable power to
obtain $R$.

\begin{theorem}[Two moving orbits are impossible: draft]\label{thm:exclusion}
Assume the three outer-trace premises stated in Section~\ref{sec:premises}.
Under the tiling hypotheses above, the noncentral tiles cannot form exactly two
$R$-orbits. In particular, using the boundary-orbit reduction below,
\[
 N=1+km\quad\Longrightarrow\quad k\geq3,
 \qquad N\geq3m+1.
\]
This does not exclude rotational Stein tilings with more orbits, or Stein
tilings without rotational symmetry.
\end{theorem}

\section{The required reductions and their dependencies}\label{sec:premises}

We use three published facts about monohedral Jordan-region disk tilings:
outer traces $S_T=T\cap\partial D$ are connected when nonempty (KLV,
Lemma~2.6); an outer trace transported to another tile is on the boundary of
its convex hull (Remark~2.7); two transported traces of positive length either
coincide or intersect only at endpoints (Remark~2.8) \cite{KLV}.
These are used for arbitrary tile congruences; empty and point traces do not
contribute length. Their statements and surrounding scope were checked in the
primary preprint, not independently reproved in this iteration.
In particular, the proof of KLV Lemma~2.6 uses tile-area equalities. Their
justification in the unrestricted positive-area Jordan-boundary case is a
specific external-review question, not settled here. This observation is not
a counterexample to the published lemma. The main draft is conditional on
the stated outer-trace premises until that review is resolved.

\paragraph{Central tile and full orbits.}
Every other tile omits $O$: otherwise its interior would accumulate at $O$
inside the open neighborhood occupied by $C$. Hence $R(C)=C$. A noncentral
tile fixed by a nonidentity power of $R$ would contain $O$ by Brouwer's fixed
point theorem applied on that Jordan disk. Thus every noncentral orbit has
size $m$.

Connectedness of $S_C$ and its rotation invariance force $S_C$ to be empty
or the whole circle. The latter would force the Jordan region $C\subset D$
to equal $D$ by Jordan separation, contradicting the other tiles. Consequently
$C$ is contained in a disk of radius $r<1$ about $O$, and
\[
 P(\conv C)\leq2\pi r<2\pi.
\]
The perimeter inequality follows, for example, by integrating convex-body
widths (the Cauchy perimeter formula); $\conv C$ has nonempty interior.

There is also an elementary proof of $S_C=\varnothing$ independent of KLV.
If $u\in C$ with $|u|=1$, then $R^ju\in C$ and
$\sum_{j=0}^{m-1}R^ju=0$. For a congruence $f(x)=Qx+t$ from $C$ to any
other tile, disk containment gives
\[
 1\geq\frac1m\sum_{j=0}^{m-1}|QR^ju+t|^2=1+|t|^2.
\]
Thus $t=f(O)=0$, putting $O$ in that other tile's interior, a contradiction.
This proves the strict central containment without boundary regularity or
outer-trace connectedness.

\paragraph{A centered matching between different outer orbits.}
The positive-length outer traces cover the circle except for finitely many
endpoints and have disjoint relative interiors. For the latter assertion,
at an interior point of an outer circular subarc a Jordan tile locally occupies
the disk side of that arc; two such tiles would overlap in their interiors.
Their total length is therefore $2\pi$. Transport all of them to $C$.
The three KLV facts and $P(\conv C)<2\pi$ force a coincidence.

For an outer representative $T$, choose $f:C\to T$ and, for $T_j=R^jT$,
choose $f_j=R^jfR^{-j}$. If $\Gamma=f^{-1}(S_T)$, its transported orbit
is $R^j\Gamma$. Its supporting unit circle has center $q=f^{-1}(O)\ne O$.
Two distinct rotated copies of $\Gamma$ cannot coincide: coincidence of their
supporting circles would give $R^{j-\ell}q=q$, hence $q=O$. Thus the forced
coincidence belongs to different outer orbits. Composing their transports
gives an isometry $h$ between two tiles and between their complete outer
arcs. Since a nondegenerate circular arc determines its supporting circle,
$h(O)=O$. There are at least two outer orbits.

\paragraph{Exactly two orbits imply alternating sectors.}
If these are the only noncentral orbits, both are outer. A centered matching
gives the same outer-arc length in both, so every outer trace has length
$\pi/m$. Around the circle the two rotational orbits alternate.
After relabeling, a centered matching between adjacent outer tiles is either
$\rho=\Rot_O(\pi/m)$ or reflection $s$ in their common endpoint radius.

If the matching is $\rho$, it swaps the two orbits and preserves their union:
$\rho^2=R$. It also preserves $C$, which is the closure of the complement of
that union in $D$. The same reasoning holds for any symmetry preserving the
noncentral tiles: the interior of each Jordan tile is dense in it, so those
tiles contain no interior point of $C$. Under the enlarged rotation $\rho$,
all noncentral tiles form one orbit, contradicting the preceding outer-orbit
argument. Therefore the matching is $s$; $R,s$ preserve the whole tiling.

Each reflection $R^js$ sends a tile in one outer orbit to a distinct tile
in the other. Its axis cannot meet that tile's interior, or the fixed point
would belong to both interiors. Connectedness of the tile interior confines
it to one component of the complement of these axes. Taking closures,
each outer tile is in a closed sector of angle $\pi/m$. Since its outer
trace has that same angle, the trace is the entire outer arc of that sector.

\paragraph{Intrinsic rotation is not an extra assumption.}
For every congruence $f:C\to T$, the tile $T$ is invariant under $fRf^{-1}$,
a rotation through $\pm2\pi/m$ about $P=f(O)\in\interior T$.
It is therefore invariant under rotation through $2\pi/m$ and its inverse.
The following lemma applies to that tile, regardless of reflected charts.

\section{The new sector calculation}

Put $\beta=\pi/(2m)$, $c=\cos\beta$, $s=\sin\beta$, and normalize the
sector bisector to the positive horizontal axis. Write
\[
 W=\{(x,y):x\geq0,\ |y|\leq x\tan\beta\},\qquad
 u_+=(c,s),\quad u_-=(c,-s).
\]
Thus $W$ has angle $\pi/m$.

\begin{lemma}[Sector obstruction, order at least three]\label{lem:sector}
Let $m\geq3$. There is no compact set $T\subset D\cap W$ with
$u_+,u_-\in T$, $O\notin T$, and rotational invariance of order $m$
about a point $P\in\interior T$.
\end{lemma}

\begin{proof}
First prove $P=(p,0)$ with $p>0$.

If $m$ is even, $T$ is centrally symmetric about $P$. All its vertical
coordinates lie in $[-s,s]$, and it contains both extrema $u_+,u_-$.
Their midpoint height is therefore $P_y=0$. The half-turn image of $u_+$
has height $-s$; the only point of $D\cap W$ at that height is $u_-$.
Consequently $P=(c,0)$ in the even case.

If $m$ is odd, let $K=\conv T$ and let
$n_+=(-s,c)$, $n_-=(-s,-c)$ be the two outward unit normals of the sector.
With $h_K(n)=\max_{x\in K}x\cdot n$, sector containment and the endpoints
give $h_K(n_+)=h_K(n_-)=0$. Rotation by
$\theta=(m-1)\pi/m$ is a power of the intrinsic rotation and sends $n_+$
to $n_-$. The centered support function $h_{K-P}$ has equal values there.
Thus $P\cdot(n_+-n_-)=0$, and $P_y=0$. Since $P\in\interior T\subset W$,
we have $p>0$.

In both cases rotate $u_+$ clockwise through $2\pi/m=4\beta$ about $P$.
The resulting point $v$ belongs to $T\subset D$. Direct expansion gives
\begin{equation}\label{eq:disk}
 |v|^2-1
 =2p(1-\cos4\beta)\left(p-\frac1{2c}\right).
\end{equation}
Indeed the expansion first gives
$2p[(1-\cos4\beta)(p-c)+s\sin4\beta]$, and
$c-s\cot(2\beta)=1/(2c)$.
The coefficient outside the final parentheses is positive, so $p\leq1/(2c)$.
For even $m\geq4$, we already have $p=c$, and $c^2>1/2$ contradicts this.

For odd $m\geq3$, rotate $u_+$ counterclockwise through
$\theta=\pi-2\beta=(m-1)\pi/m$. The result is
\begin{equation}\label{eq:ray}
 z=(2pc-1)(c,-s)\in T.
\end{equation}
Sector containment implies $2pc-1\geq0$. Combined with the disk bound it
forces $p=1/(2c)$, hence $z=O\in T$, a contradiction.
\end{proof}

Only the endpoints were needed in this lemma, not the full outer arc, Jordan
regularity, connected contacts, or any equations on the remaining contour.

\begin{lemma}[The remaining order-two lens]\label{lem:lens}
Let $m=2$ and let a compact set $T\subset D\cap W$ contain $u_+,u_-$ and
be centrally symmetric about $P$. Then
\[
 P=(1/\sqrt2,0),\qquad
 T\subset D\cap(2P-D),\qquad
 \area(T)\leq\frac\pi2-1.
\]
\end{lemma}

\begin{proof}
The even-case endpoint argument above gives $P=(c,0)$, now $c=1/\sqrt2$.
Central symmetry and $T\subset D$ give $T\subset2P-D$. The two unit disks
have centers separated by $2c=\sqrt2$. Their intersection is a lens made of
two circular segments of opening angle $\pi/2$, each of area $\pi/4-1/2$.
This proves the bound.
\end{proof}

\section{Finishing the proposed obstruction}

For $m\geq3$, the reductions produce an outer tile satisfying
Lemma~\ref{lem:sector}, a contradiction. For $m=2$ there are five congruent
tiles. Lemma~\ref{lem:lens} bounds the area of every tile, including $C$.
Subadditivity of Lebesgue measure would give
\[
 \pi=\area(D)\leq\sum_T\area(T)
 \leq5\left(\frac\pi2-1\right)<\pi,
\]
where the strict inequality is equivalent to $\pi<10/3$, for example a
consequence of $\pi<22/7$. This proves Theorem~\ref{thm:exclusion}, subject
to the stated proof-review status and external premises.

Importantly, this uses subadditivity, not $\area(T)=\pi/5$. Arbitrary Jordan
curves may have positive area; disjoint interiors alone do not justify adding
tile areas as equal shares. No assumption about zero-area tile boundaries is
needed here.

\section{Limits, attempted falsification and use}

The even/odd split uses only actual intrinsic rotations of a geometric tile,
not graph automorphisms. The sector argument has not been asserted for a
shorter outer trace or a wider sector. Order two is not rejected by the
sector lemma: its lens is a genuine centrally symmetric shape fitting the
quarter-disk, so the separate coverage-area step is essential. No point-only
contact has been promoted to a positive outer arc.

For another assumption check, when $m=3$ the equilateral triangle with vertices
$O,u_+,u_-$ fits the sector and the disk and is intrinsically threefold
symmetric. It violates precisely the requirement $O\notin T$ of the sector
lemma. Thus that requirement must not be dropped or silently replaced by
``$O$ is not an interior point of $T$''.

The proposed count reduction must not be enabled in code until independent
proof review and a separate premise-detector audit. Nothing here establishes
nonexistence of all Stein tilings or of known weaker nontrivial disk tilings.
A limited prior-art search did not certify novelty. An external review
contract is recorded under \texttt{docs/research/contracts/}.

\begin{thebibliography}{1}
\bibitem{KLV}
A. Kurusa, Z. L\'angi and V. V\'igh,
\emph{Tiling a Circular Disc with Congruent Pieces},
Mediterranean Journal of Mathematics 17 (2020), Article 156.
\href{https://doi.org/10.1007/s00009-020-01595-3}{doi:10.1007/s00009-020-01595-3}.
Statement-level checks used the primary preprint
\href{https://arxiv.org/html/1910.03836}{arXiv:1910.03836v1},
Lemma 2.6 and Remarks 2.7--2.8 (accessed 21 October 2026).
\end{thebibliography}
\end{document}
